\subsection{Đề 5}
\Opensolutionfile{ans}[ans/ans-CTST_1-OTGK1_De5-LC]
\caulc
\begin{ex}%[Câu 1]%[1D1N1-2]
	Góc có số đo $45^{\circ}$ bằng bao nhiêu radian
	\choice
	{\True $\dfrac{\pi}{4}$}
	{$\dfrac{2\pi}{3}$}
	{$\dfrac{\pi}{3}$}
	{$\dfrac{\pi}{2}$}
	\loigiai{
		Ta có $45^{\circ} = \dfrac{45 \cdot \pi}{180} = \dfrac{\pi}{4}$ rad.
	}
\end{ex}
\begin{ex}%[Câu 2]%[1D1N2-2]
	Cho $\cos \alpha = \dfrac{5}{13}$ với $-\dfrac{\pi}{2} < \alpha < 0$. Giá trị của $\sin \alpha$ là
	\choice
	{$\dfrac{12}{13}$}
	{\True $-\dfrac{12}{13}$}
	{$-\dfrac{5}{13}$}
	{$\dfrac{13}{5}$}
	\loigiai{
		Ta có $\sin^2\alpha + \cos^2\alpha = 1 \Leftrightarrow \sin^2\alpha = 1 - \cos^2\alpha = 1 - \left(\dfrac{5}{13}\right)^2 = \left(\dfrac{12}{13}\right)^2$.\\
		$\Leftrightarrow \sin\alpha = \dfrac{12}{13}$ hoặc $\sin\alpha = -\dfrac{12}{13}$.\\
		Vì $-\dfrac{\pi}{2} < \alpha < 0$ nên $\sin\alpha < 0$.\\
		Vậy $\sin\alpha = -\dfrac{12}{13}$.
	}
\end{ex}
\begin{ex}%[Câu 3]%[1D1N3-2]
	Giá trị lượng giác $\sin \left(\dfrac{5\pi}{12}\right)$ bằng
	\choice
	{$0{,}9$}
	{$\dfrac{\sqrt{2} \left(1 + \sqrt{3}\right)}{2}$}
	{$\dfrac{\sqrt{3} \left(1 + \sqrt{2}\right)}{4}$}
	{\True $\dfrac{\sqrt{2} \left(1 + \sqrt{3}\right)}{4}$}
	\loigiai{
		Ta có 
		\begin{align*}
			\sin \left(\dfrac{5\pi}{12}\right) &= \sin \left(\dfrac{2\pi + 3\pi}{12}\right) = \sin \left(\dfrac{\pi}{6} + \dfrac{\pi}{4}\right) = \sin \left(\dfrac{\pi}{6}\right) \cos \left(\dfrac{\pi}{4}\right) + \cos \left(\dfrac{\pi}{6}\right) \sin \left(\dfrac{\pi}{4}\right)\\
			&=\dfrac{1}{2}\cdot \dfrac{\sqrt{2}}{2} + \dfrac{\sqrt{3}}{2} \cdot \dfrac{\sqrt{2}}{2} = \dfrac{\sqrt{2} \left(1 + \sqrt{3}\right)}{4}.
		\end{align*}
	}
\end{ex}
\begin{ex}%[Câu 4]%[1D1N4-4]
	Hàm số nào sau đây là hàm số lẻ trên tập xác định của nó?
	\choice
	{\True $y=\sin x$}
	{$y=\cos x$}
	{$y=\tan^2 x$}
	{$y=x+1$}
	\loigiai{
		\begin{itemize}
			\item $\sin(-x) = - \sin x$ nên $y=\sin x$ là hàm số lẻ.
			\item $\cos(-x) = \cos(x)$ nên hàm số $y=\cos x$ là hàm số chẵn.
			\item $\tan^2 (-x) = \left[\tan(-x)\right]^2 = \left(-\tan x\right)^2 = \tan^2 x$ nên hàm số $y=\tan^2 x$ là hàm số chẵn.
			\item $(-x) + 1 \ne (x) + 1$ và $(-x) + 1 \ne -(x + 1)$ nên hàm số $y=x+1$ không có tính chẵn, lẻ.
		\end{itemize}
	}
\end{ex}
\begin{ex}%[Câu 5]%[1D1N5-3]
	Nghiệm của phương trình $\sin x = 0$ là
	\choice
	{$x=k2\pi$, $k \in \mathbb{Z}$}
	{\True $x=k\pi$, $k \in \mathbb{Z}$}
	{$x=\dfrac{\pi}{2} + k\pi$, $k \in \mathbb{Z}$}
	{$x=\dfrac{\pi}{2} + k2\pi$, $k \in \mathbb{Z}$}
	\loigiai{
		Ta có $\sin x = 0 \Leftrightarrow x = k\pi$, $k \in \mathbb{Z}$.
	}
\end{ex}
\begin{ex}%[Câu 6]%[1D2N1-3]
	Số hạng thứ $3$ của dãy số $\heva{&u_1 = 1\\&u_{n+1} = \dfrac{u_n}{1+u_n} \quad (n \ge 1)}$ là	
	\choice
	{$u_3=\dfrac{1}{2}$}
	{$u_3=\dfrac{1}{4}$}
	{\True $u_3=\dfrac{1}{3}$}
	{$u_3=1$}
	\loigiai{
		Dãy số $\heva{&u_1 = 1\\&u_{n+1} = \dfrac{u_n}{1+u_n} \quad (n \ge 1)}$ có $u_2 = \dfrac{u_1}{1+u_1} \Rightarrow u_2 = \dfrac{1}{1+1}  \Rightarrow u_2 = \dfrac{1}{2}$ và $u_3 = \dfrac{1}{3}$.
	}
\end{ex}
\begin{ex}%[Câu 7]%[1D2N2-1]
	Dãy số nào sau đây là cấp số cộng?
	\choice
	{\True $3$; $5$; $7$; $9$; $11$}
	{$3$; $5$; $7$; $9$; $12$}
	{$3$; $5$; $7$; $9$; $10$}
	{$3$; $5$; $6$; $7$; $9$}
	\loigiai{
		Xét dãy số $3$; $5$; $7$; $9$; $11$ có
		\begin{itemize}
			\item $5=3 + 2$.
			\item $7=5 + 2$.
			\item $9=7 + 2$.
			\item $11=9 + 2$.
		\end{itemize}
		Đây là cấp số cộng với $u_1=3$ công sai $d=2$.
	}
\end{ex}
\begin{ex}%[Câu 8]%[1D2N3-2]
	Cho cấp số nhân $32$; $16$; $8$; $4$; $2$ có công bội là
	\choice
	{$q=2$}
	{\True $q=\dfrac{1}{2}$}
	{$q=\dfrac{1}{4}$}
	{$q=\dfrac{1}{3}$}
	\loigiai{
		Cho cấp số nhân $32$; $16$; $8$; $4$; $2$ có công bội là $q=\dfrac{1}{2}$.
	}
\end{ex}
\begin{ex}%[Câu 9]%[1H4N1-2]
	Cho hình chóp $S.ABCD$, gọi $O$ là giao điểm của $AC$ và $BD$. Lấy $M$, $N$ lần lượt thuộc các cạnh $SA$, $SC$. Đường thẳng nào sau đây \textbf{không} thuộc mặt phẳng $(SAC)$.
	\begin{center}
		\begin{tikzpicture}[scale=0.7, font=\footnotesize,line join=round, line cap=round, >=stealth]
			\path 
			(0,0) coordinate (A)
			++(-140:3) coordinate (B)
			++(0:5) coordinate (C)
			($(A)+(C)-(B)$) coordinate (D)
			($(A)+(1,4)$) coordinate (S)
			($(S)!0.5!(A)$) coordinate (M)
			($(S)!0.5!(C)$) coordinate (N)
			($(A)!0.5!(C)$) coordinate (O)
			;
			\foreach \i in{B,C,D}{\draw (S)--(\i);};
			\draw (B)--(C)--(D);
			\draw[dashed] (S)--(A)--(B)--(D)--(A)--(C) (M)--(N);
			\foreach \i/\g in {A/-90,B/-90,C/-90,D/-60,S/90,M/30,N/30,O/-90}
			\fill[black] (\i) circle(1pt)+(\g:5mm)node[scale=1]{$\i$};
		\end{tikzpicture}	
	\end{center}
	\choice
	{$SA$}
	{$MN$}
	{$AC$}
	{\True $BO$}
	\loigiai{
		\begin{center}
			\begin{tikzpicture}[scale=0.7, font=\footnotesize,line join=round, line cap=round, >=stealth]
				\path 
				(0,0) coordinate (A)
				++(-140:3) coordinate (B)
				++(0:5) coordinate (C)
				($(A)+(C)-(B)$) coordinate (D)
				($(A)+(1,4)$) coordinate (S)
				($(S)!0.5!(A)$) coordinate (M)
				($(S)!0.5!(C)$) coordinate (N)
				($(A)!0.5!(C)$) coordinate (O)
				;
				\foreach \i in{B,C,D}{\draw (S)--(\i);};
				\draw (B)--(C)--(D);
				\draw[dashed] (S)--(A)--(B)--(D)--(A)--(C) (M)--(N);
				\foreach \i/\g in {A/-90,B/-90,C/-90,D/-60,S/90,M/30,N/30,O/-90}
				\fill[black] (\i) circle(1pt)+(\g:5mm)node[scale=1]{$\i$};
			\end{tikzpicture}	
		\end{center}
		Ta có $S$, $A$, $C$, $M$, $N \in (SAC)$ nên $SA$, $SC$, $MN\subset (SAC)$.\\
		Mặt khác $\heva{&B \not\in (SAC)\\&O\in (SAC)\\} \Rightarrow BO \not\subset (SAC)$.
	}
\end{ex}
\begin{ex}%[Câu 10]%[1H4N2-2]
	Cho hình hộp chữ nhật $ABCD.A'B'C'D'$ như hình. Mệnh đề nào sau đây là \textbf{sai}?
	\begin{center}
		\begin{tikzpicture}[scale=0.6, font=\footnotesize,line join=round, line cap=round, >=stealth]
			\path 
			(0,0) coordinate (A)
			++(-130:3) coordinate (B)
			++(0:6) coordinate (C)
			($(A)+(C)-(B)$) coordinate (D)
			($(A)!1/2!(C)$) coordinate (O)
			;
			\foreach \i in {A,B,C,D}{
				\coordinate (\i') at ($(\i)+(0,4)$);
			}
			\draw (A')--(B')--(C')--(D')--cycle;
			\draw (B)--(B') (C)--(C') (D)--(D')  (B)--(C)--(D) (A')--(C');
			\draw[dashed,thin](B)--(A)--(A') (C)--(A)--(D);
			\foreach \i/\g in {A'/90,B'/90,C'/90,D'/90,A/-90,B/-90,C/-90,D/-90}
			\fill[black] (\i) circle(1pt)+(\g:5mm)node[scale=1]{$\i$};
		\end{tikzpicture}
	\end{center}
	\choice
	{$AB \parallel CD$}
	{$AB \parallel C'D'$}
	{$AB \parallel A'B'$}
	{\True $AB \parallel A'C'$}
	\loigiai{
		Rõ ràng là $AB$ không song song với $A'C'$.
	}
\end{ex}
\begin{ex}%[Câu 11]%[1D1H5-5]
	Nghiệm của phương trình $\cos \left(\dfrac{x}{2}\right)= -\dfrac{1}{2}$ là
	\choice
	{\True $x=\dfrac{4\pi}{3}+k2\pi$ hoặc $x=-\dfrac{4\pi}{3} + k2\pi, \ k\in \mathbb{Z}$}
	{$x=\dfrac{2\pi}{3} + k2\pi$ hoặc $x=-\dfrac{2\pi}{3} + k2\pi, \ k\in \mathbb{Z}$}
	{$x=\dfrac{4\pi}{3} + k\pi$ hoặc $x=-\dfrac{4\pi}{3} + k\pi, \ k\in \mathbb{Z}$}
	{$x=\dfrac{2\pi}{6} + k\pi$ hoặc $x=-\dfrac{2\pi}{6} + k\pi, \ k\in \mathbb{Z}$}
	\loigiai{
		Ta có 
		$-\dfrac{1}{2} = -\cos \left(\dfrac{\pi}{3}\right) = \cos \left(\pi - \dfrac{\pi}{3}\right) = \cos \left(\dfrac{2\pi}{3}\right)$.\\
		Khi đó \\
		\begin{align*}
			&\cos \left(\dfrac{x}{2}\right) = - \dfrac{1}{2}.\\
			\Leftrightarrow &\cos \left(\dfrac{x}{2}\right) = \cos \left(\dfrac{2\pi}{3}\right).\\
			\Leftrightarrow&\hoac{&\dfrac{x}{2} = \dfrac{2\pi}{3} + k2\pi\\&\dfrac{x}{2}=-\dfrac{2\pi}{3} + k2\pi} \\
			\Leftrightarrow &\hoac{&x = \dfrac{4\pi}{3} + k2\pi\\&x=-\dfrac{4\pi}{3} + k2\pi} \quad k\in \mathbb{Z}.
		\end{align*}
	}
\end{ex}
\begin{ex}%[Câu 12]%[1D2H2-4]
	Cho cấp số cộng $(u_n)$ có số hạng đầu $u_1 = -2$ và công sai $d = 3$. Số $-168$ là số hạng thứ bao nhiêu?
	\choice
	{$28$}
	{$30$}
	{\True $29$}
	{$31$}
	\loigiai{
		Ta có \\
		\begin{align*}
			& u_n = u_1(n-1)d \\
			\Leftrightarrow & -168=-2\cdot (n-1)\cdot 3.\\ 
			\Leftrightarrow & n = \dfrac{-168}{-6} + 1 = 29.
		\end{align*}
	}
\end{ex}
\Closesolutionfile{ans}
\indapan{6}{ans/ans-CTST_1-OTGK1_De5-LC}
%-------------------------------------------
\Opensolutionfile{ans}[ans/ans-CTST_1-OTGK1_De5-DS]
\cauds
\begin{ex}%[Câu 1]%[1D1N5-5]
	Cho hàm số $y = \sin x$. Khi đó
	\choiceTF
	{\True $\sin x < 0 $ khi $ -\dfrac{\pi}{2} < x < 0$}
	{\True Hàm số $y = \sin x$ là hàm số lẻ với mọi $x \in \mathbb{R}$}
	{Phương trình $\sin x = 1$ có nghiệm $x=\dfrac{\pi}{2} + k\pi,\ k\in \mathbb{Z}$}
	{Hàm số $y=\sin x$ có chặn dưới là $0$}
	\loigiai{
		\begin{itemchoice}
			\itemch Đúng. Vì $x$  thuộc gốc phần tư thứ tư nên $\sin x < 0 $.
			\itemch Đúng. Vì $f(-x) = \sin\left(x\right) = - \sin x = -f(x)$.
			\itemch Sai.\\
				\begin{align*}
					&\sin x = 1\\
					\Leftrightarrow &x = \dfrac{\pi}{2} + k\mathbf{2}\pi\ \forall k \in \mathbb{Z}.
				\end{align*}
			\itemch Sai. Vì $-1 \le \sin x \le 1$ nên hàm số có chặn dưới là $-1$.
		\end{itemchoice}
	}
\end{ex}
\begin{ex}%[Câu 2]%[1D1H2-2]
	Cho $\sin \alpha = \dfrac{12}{13}$ và $\dfrac{\pi}{2}<\alpha<\pi$. Khi đó
	\choiceTF
	{\True $\cos\alpha = -\dfrac{5}{13}$}
	{$\cos\alpha = \pm \dfrac{5}{13}$}
	{\True $\tan\alpha = -\dfrac{12}{5}$}
	{$\cot\alpha = \dfrac{5}{13}$}
	\loigiai{
		Với $\sin \alpha = \dfrac{12}{13}$ và $\dfrac{\pi}{2}<\alpha<\pi$.\\
		Ta có \\
		$\sin^2\alpha+\cos^2\alpha=1 \Leftrightarrow \cos\alpha = \pm\dfrac{5}{13} \Rightarrow \cos\alpha = -\dfrac{5}{13}\ \left(\dfrac{\pi}{2}<\alpha<\pi\right)$.\\
		$\Rightarrow \tan\alpha = \dfrac{\sin\alpha}{\cos\alpha} = -\dfrac{12}{5}$.\\
		$\Rightarrow \cot\alpha=\dfrac{1}{\tan\alpha}=-\dfrac{12}{5}$.
	}
\end{ex}
\begin{ex}%[Câu 3]%[1D2H1-5]
	Cho dãy số $(u_n)$ được xác định bởi $\heva{&u_1=3\\&u_{n+1}=2u_n \quad (n \ge 1)}$. Khi đó
	\choiceTF
	{Dãy số $(u_n)$ là dãy số giảm}
	{Dãy số $(u_n)$ là dãy số bị chặn}
	{\True $u_2 = 6$}
	{\True Công thức tổng quát của $(u_n)$ là $u_n=2^{n-1}\cdot 3$}
	\loigiai{
		\begin{itemchoice}
			\itemch Sai. Dãy số $(u_n)$ là dãy số tăng. Vì $u_n < u_{n+1},\ \forall n \in \mathbb{N}^*$.
			\itemch Sai. Vì khi $n$ tiến ra $+\infty$ thì $u_n = \left(2^{n-1}\cdot 3\right)$ cũng tiến ra $+\infty\ \forall n \in \mathbb{N}^{\ast}$. Do đó dãy số $(u_n)$ không có chặn trên. Nên $(u_n)$ không bị chặn.
			\itemch Đúng. Vì $u_2 = 2u_1=2\cdot 3=6$.
			\itemch Đúng. \\
				\textbf{Cách 1: Dùng quy nạp.}\\
				Ta có $u_1 = 3$.\\
				Giả sử $u_k = 2^{k-1}\cdot 3$, đúng $\forall k \ge 1$.\\
				Ta chứng minh $u_{k+1} = 2^{k}\cdot 3$.
				Thật vậy $u_{k+1} = 2u_k \Rightarrow 2u_k = 2^k \cdot 3 \Rightarrow u_k = 2^{k-1}\cdot 3$.\\
				Vậy $u_n = 2^{n-1}\cdot 3$, $\forall n \ge 1$.\\
				\textbf{Cách 2: Không dùng quy nạp.}\\
				Dự đoán công thức tổng quát của dãy $(u_n)$ là $u_n = 2^{n-1}\cdot 3$.\\
				Ta kiểm tra lại: 
				\begin{itemize}
					\item $u_1 = 2^{1-1}\cdot 3 = 2^0\cdot 3 = 3$.
					\item Thay $n = k$ ta được $u_k = 2^{k-1}\cdot 3$. Thay $k = n+1$ ta được $u_{n+1} = u^n \cdot 3 = 2u^{n-1}\cdot 3 = 2u_n$.
				\end{itemize}			
		\end{itemchoice}
	}
\end{ex}
\begin{ex}%[Câu 4]%[1H4H3-2]
	\immini[thm]{Cho hình chóp $S.ABCD$, như hình vẽ, có đáy $ABCD$ là hình bình hành. Gọi $O$ giao điểm của $AC$ và $BD$; điểm $M$, $N$ lần lượt là trung điểm của $SA$ và $SC$. Khi đó
	\choiceTF
	{\True Gọi $P$ là giao điểm của $SO$ và $MN$. Khi đó điểm $P$ nằm trên mặt phẳng $\left(SBD\right)$}
	{\True Đường thẳng $AC$ song song với $\left(DMN\right)$}
	{Giao tuyến của $(SAD)$ và $(SBC)$ sẽ song song với đường thẳng $AB$}
	{\True Đường thẳng $MN$ song song với mặt phẳng $(ABCD)$}
}
{
	\begin{tikzpicture}[scale=0.7, font=\footnotesize,line join=round, line cap=round, >=stealth]
		\path 
		(0,0) coordinate (A)
		++(-140:3) coordinate (B)
		++(0:6) coordinate (C)
		($(A)+(C)-(B)$) coordinate (D)
		($(A)+(1,3)$) coordinate (S)
		($(A)!0.5!(C)$) coordinate (O)
		($(S)!0.5!(A)$) coordinate (M)
		($(S)!0.5!(C)$) coordinate (N)
		($(M)!0.5!(N)$) coordinate (P)
		;
		\foreach \i in{B,C,D}{\draw (S)--(\i);};
		\draw (B)--(C)--(D);
		\draw[dashed] (S)--(A)--(B) (A)--(D) (A)--(C) (B)--(D) (S)--(O) (M)--(N);
		\foreach \i/\g in {A/-90,B/-90,C/-90,D/0,S/90,M/30,N/0,O/-90, P/-150}
		\fill[black] (\i) circle(1pt)+(\g:4mm)node[scale=1]{$\i$};
	\end{tikzpicture}
}
	\loigiai{
		\begin{center}
			\begin{tikzpicture}[scale=0.7, font=\footnotesize,line join=round, line cap=round, >=stealth]
				\path 
				(0,0) coordinate (A)
				++(-140:3) coordinate (B)
				++(0:6) coordinate (C)
				($(A)+(C)-(B)$) coordinate (D)
				($(A)+(1,3)$) coordinate (S)
				($(S)-(1,0)$) coordinate (Sx)
				($(S)+(4,0)$) coordinate (Sxx)
				($(A)!0.5!(C)$) coordinate (O)
				($(S)!0.5!(A)$) coordinate (M)
				($(S)!0.5!(C)$) coordinate (N)
				($(M)!0.5!(N)$) coordinate (P)
				;
				\foreach \i in{B,C,D}{\draw (S)--(\i);};
				\draw (B)--(C)--(D) (Sx)--(Sxx);
				\draw[dashed] (S)--(A)--(B) (A)--(D) (A)--(C) (B)--(D) (S)--(O) (M)--(N);
				\foreach \i/\g in {A/-90,B/-90,C/-90,D/0,S/90,M/30,N/0,O/-90, P/-150}
				\fill[black] (\i) circle(1pt)+(\g:4mm)node[scale=1]{$\i$};
				\fill[black] (Sxx) circle(0pt) + (90:4mm) node[scale=1]{$x$};
			\end{tikzpicture}
		\end{center}
		\begin{itemchoice}
			\itemch Đúng. $P=MN\cap SO \Rightarrow P \in SO \subset (SBD) \Rightarrow P \in (SBD)$.
			\itemch Đúng. Ta có $AC \parallel MN \subset (DMN)$ nên $AC \parallel (DMN)$.
			\itemch Sai. $(SAD) \cap (SBC) = Sx \parallel AD \parallel BC$ nên $Sx$ không song song với $AB$.
			\itemch Đúng. $MN$ là đường trung bình của tam giác $SAC$ nên $MN \parallel AC \subset (ABCD)$. Suy ra $MN \parallel (ABCD)$.
		\end{itemchoice}
	}
\end{ex}
\Closesolutionfile{ans}
\indapan{3}{ans/ans-CTST_1-OTGK1_De5-DS}
%-------------------------------------------
\Opensolutionfile{ans}[ans/ans-CTST_1-OTGK1_De5-KQ]
\caukq
\begin{ex}%[Câu 1]%[1D1V4-8]
	Khi đu quay hoạt động, vận tốc theo phương ngang của một cabin $M$ phụ thuộc vào góc lượng giác $\alpha = (Ox, OM)$ theo hàm số $v_x = 0{,}25\sin\alpha$ (m/s). Vận tốc lớn nhất của cabin là
	\begin{center}
		\begin{tikzpicture}[scale=0.7,font=\small,line join=round,line cap=round,>=stealth]
			\tikzset{declare function={g=0;}}
			\coordinate (O) at (0,0);
			%	\coordinate[label=above right:$A$] (A) at (0:4.4);
			%	\coordinate[label=above right:$M$] (M) at (30+g:4.4);
			\pgfmathsetmacro{\gx}{60-g};
			\coordinate (v) at ($(M)+(120+g:3.3)$);
			\coordinate (vx) at ($(M)+(180:3.3*cos \gx)$);
			\draw[thick] (O) circle(4.4) (O) circle(3.53);
			\foreach \t in {0,1,...,11}{
				\draw[very thick,fill] (O)--(g+30*\t:4.4)circle(0.06);
				\begin{scope}[shift={(g+30*\t:4.4)}]
					\draw[fill=teal] (-0.27,0)--(0.27,0)--(0.36,-0.18)--(-0.36,-0.18)--cycle
					(0.32,-0.18)--(0.32,-0.54) (-0.32,-0.18)--(-0.32,-0.54)
					(-0.4,-0.54)--(0.4,-0.54)--(0.27,-1)--(-0.27,-1)--cycle;
			\end{scope}}
			\draw[fill=gray!70!yellow] (0.2,0)--(2.35,-5.53)--(3.65,-5.53)--(3.65,-6.02)--(1.95,-6.02)--(0,-0.85)--(-1.95,-6.02)--(-3.65,-6.02)--(-3.65,-5.53)--(-2.35,-5.53)--(-0.2,0)--cycle;
			\draw[fill=gray!10!yellow] (O) circle(0.34);
			\draw[fill] (-5,-5.85)--(0,-5.85) circle(0.08)--(6.4,-5.85);
			\draw[->] (-5,0)--(6.7,0) node [below]{\color{black}$x$};
			\path 
			(3.82,2.22) coordinate (M)
			($(M)-(1.8,0)$) coordinate (N)
			($(N)+(0,2.5)$) coordinate (P)
			;
			\fill[black] (M) circle(1pt) node[above right]{$M$};
			\fill[black] (N) circle(0pt) node[below left]{$\overrightarrow{v}_x$};
			\fill[black] (P) circle(0pt) node[above]{$\overrightarrow{v}$};
			\fill[black] (O) circle(2pt);
			\fill[black] (O)+(-90:0.3) circle(0pt) node[below]{$O$};
			
			\draw[->] (M)--(N);
			\draw[->] (M)--(P);
			\draw[dashed] (N)--(P);
			\draw (O)--(M);
			\draw[->] (O)+(0:0.8) arc (0:30:0.8) node[below right]{$\alpha$};
			\draw pic[draw, angle radius = 4mm] {right angle=O--M--P};
		\end{tikzpicture}
	\end{center}
	\shortans{$0{,}25$}
	\loigiai{
		Với mọi  Ta có  $-1\le \sin\alpha \le 1$, $\forall \alpha\in \mathbb{R}$.\\
		Do đó $-0{,}25 \le 0{,}25\sin\alpha \le 0{,}25 \Leftrightarrow -0{,}25 \le v(x) \le 0{,}25$.\\
		Vậy vận tốc lớn nhất của cabin là $0{,}25$.
	}
\end{ex}
\begin{ex}%[Câu 2]%[1D1V5-6]
	Một chiếc lò xo một đầu cố định, một đầu gắn chất điểm $A$. Nén lò xo một lực làm chất điểm $A$ dao động xung quanh vị $O$, quãng đường chất điểm $A$ di chuyển được xác định bởi công thức $s=2\sin\left(2t+\dfrac{\pi}{4}\right)$ (cm), với $t$ là thời gian được tính bằng giây~(s). Trong $3$ giây đầu, thời điểm nào thì $s = \sqrt{3}$ (cm) (Làm tròn kết quả đến hàng phần trăm)?
	\shortans{$0{,}13$}
	\loigiai{
		Ta có\\
		\begin{align*}
			&2\sin\left(2t+\dfrac{\pi}{4}\right) = \sqrt{3} \\
			\Leftrightarrow &\sin \left(2t+\dfrac{\pi}{4}\right) = \dfrac{\sqrt{3}}{2} = \sin\left(\dfrac{\pi}{3}\right)\\ \Leftrightarrow &\hoac{&2t + \dfrac{\pi}{4} = \dfrac{\pi}{3} + k2\pi\\&2t + \dfrac{\pi}{4} = \pi - \dfrac{\pi}{3} + k2\pi\\} \\
			\Leftrightarrow &\hoac{&t = \dfrac{\pi}{24} + k\pi\\&t=\dfrac{5\pi}{24} + k\pi}\quad k \in \mathbb{Z}.\\
		\end{align*}
		Với $0<t\le 3$, ta chọn $t=\dfrac{\pi}{24}\approx 0{,}13$ với $k=0$.
	}
\end{ex}
\begin{ex}%[Câu 3]%[1D2V2-7]
	Khán đài D của một sân vận động có $20$ hàng ghế xếp theo hình quạt. Hàng thứ nhất có $13$ ghế, hàng thứ hai có $16$ ghế, hàng thứ ba có $19$ ghế, $\ldots$, cứ thế tiếp tục cho đến hàng cuối cùng. Số ghế ở hàng cuối cùng là
	\shortans{$741$}
	\loigiai{
		Đây là cấp số cộng với $u_1 = 13$, $d = 3$.\\ Khi đó $u_{20} = u_1(20-1)d=13\cdot 19 \cdot 3 = 741$ ghế.
	}
\end{ex}
\begin{ex}%[Câu 4]%[1D2V3-8]
	Tỉnh Trà Vinh năm $2022$ có dân số là $1\,019\,258$ người và tốc độ gia tăng dân số trung bình mỗi năm là $0{,}06$ \%. Dự đoán dân số của tỉnh Trà Vinh năm $2030$ là bao nhiêu (Làm tròn kết quả đến hàng phần nghìn)?
	\shortans{$1\,024$}
	\loigiai{
		Đây là cấp số nhân với $u_1 = 1\,019\,258$, công bội $q=1{,}06$. \\
		Năm $2030$ là năm thứ $9$, tức là dân số tỉnh Trà Vinh năm $2030$ dự đoán là\\ \centerline{$u_9 = u_1\cdot q^{9-1} = 1\,019\,258 \cdot \left(1+\dfrac{0{,}06}{100}\right)^{8} = 1\,024\,160{,}725\approx 1\,024$ nghìn người.}
	}
\end{ex}
\begin{ex}%[Câu 5]%[1H4V2-6]
	Cho hình chóp $S.ABCD$ có đáy là hình bình hành, cạnh $AB=6$ cm, $AC$ và $BD$ cắt nhau tại $O$. Gọi $I$ là trung điểm của $SO$. Mặt phẳng $(ICD)$ cắt $SA$, $SB$ lần lượt tại $M$, $N$. Độ dài $MN$ là
	\begin{center}
		\begin{tikzpicture}[scale=0.5, font=\footnotesize,line join=round, line cap=round, >=stealth]
			\path 
			(0,0) coordinate (A)
			++(-140:3) coordinate (B)
			++(0:8) coordinate (C)
			($(A)+(C)-(B)$) coordinate (D)
			($(A)+(1,6)$) coordinate (S)
			($(A)!0.5!(C)$) coordinate (O)
			($(S)!0.5!(O)$) coordinate (I)
			;
			\foreach \i in{B,C,D}{\draw (S)--(\i);};
			\draw (B)--(C)--(D);
			\draw[dashed] (O)--(S)--(A)--(B)--(D)--(A)--(C) (D)--(I)--(C);
			\foreach \i/\g in {A/-90,B/-90,C/-90,D/-90,S/90,O/-90,I/30}
			\fill[black] (\i) circle(1pt)+(\g:4mm)node[scale=1]{$\i$};
		\end{tikzpicture}
	\end{center}
	\shortans{$2$}
	\loigiai{
		\begin{center}
			\begin{tikzpicture}[scale=0.5, font=\footnotesize,line join=round, line cap=round, >=stealth]
				\path 
				(0,0) coordinate (A)
				++(-140:3) coordinate (B)
				++(0:8) coordinate (C)
				($(A)+(C)-(B)$) coordinate (D)
				($(A)+(1,6)$) coordinate (S)
				($(A)!0.5!(C)$) coordinate (O)
				($(S)!0.5!(O)$) coordinate (I)
				($(S)!1/3!(A)$) coordinate (M)
				($(S)!1/3!(B)$) coordinate (N)
				;
				\foreach \i in{B,C,D}{\draw (S)--(\i);};
				\draw (B)--(C)--(D);
				\draw[dashed] (O)--(S)--(A)--(B)--(D)--(A)--(C) (D)--(I)--(C) (I)--(M)--(N)--(I);
				\foreach \i/\g in {A/-90,B/-90,C/-90,D/-90,S/90,O/-90,I/30,N/150}
				\fill[black] (\i) circle(1pt)+(\g:4mm)node[scale=1]{$\i$};
				\fill[black] (M) circle(1pt)+(0:4mm)node[scale=1]{$M$};
			\end{tikzpicture}
		\end{center}
		Ta có $\heva{&I \in (ICD)\\&I \in SO \subset (SAC)\\&C \in (ICD)\\&C \in (SAC)} \Rightarrow IC$ là giao tuyến của $(ICD)$ và $(SAC)$.\\
		Mặt khác điểm $M$ là giao điểm của $SA$ và $(ICD)$, mà $SA \subset (SAC)$ nên $M$ nằm trên giao tuyến của $(ICD)$ và $(SAC)$.\\
		Suy ra $M$, $I$, $C$ thẳng hàng.
		Tương tự, ta cũng chứng minh được $N$, $I$, $D$ thẳng hàng.\\
		Áp dụng Định lý Menelaus trong $\triangle SAO$\\
		$$\dfrac{CA}{CO}\cdot \dfrac{IO}{IS} \cdot \dfrac{MS}{MA} = 1\Rightarrow 2\cdot 1\cdot \dfrac{MS}{MA} = 1\Rightarrow \dfrac{MS}{MA} = \dfrac{1}{2}\Rightarrow \dfrac{SM}{SA} = \dfrac{1}{3}.$$
		Ta thấy $AB$, $CD$, $MN$ lần lượt là giao tuyến của ba mặt phẳng $(SAB)$, $(SCD)$ và $(ICD)$.\\
		Theo định lý về giao tuyến của ba mặt phẳng và $AB \parallel CD$ nên $MN \parallel AB \parallel CD$.\\
		Áp dụng Định lý Thales trong $\triangle SAB$
		$$\dfrac{MN}{AB} = \dfrac{SM}{SA} = \dfrac{1}{3}\Rightarrow MN = \dfrac{1}{3}\cdot AB = \dfrac{1}{3}\cdot 6 = 2.$$
	}
\end{ex}
\begin{ex}%[Câu 6]%[1H4V3-5]
	Cho hình chóp $S.ABCD$ có đáy $ABCD$ là hình bình hành và một điểm $M$ nằm trên cạnh $AD$ (giữa $A$ và $D$) sao cho $AD = 3MD$. Một mặt phẳng $(\alpha)$ đi qua $M$, song song với $CD$ và $SA$, cắt $BC$, $SC$, $SD$ lần lượt tại $N$, $P$, $Q$. Với cạnh $CD = 9{,}3$ (cm) thì độ dài đoạn $PQ$ là bao nhiêu?
	\begin{center}
		\begin{tikzpicture}[scale=0.5, font=\footnotesize,line join=round, line cap=round, >=stealth]
			\path 
			(0,0) coordinate (A)
			++(-140:3) coordinate (B)
			++(0:8) coordinate (C)
			($(A)+(C)-(B)$) coordinate (D)
			($(A)+(1,6)$) coordinate (S)
			($(D)!1/3!(A)$) coordinate (M)
			($(C)!1/3!(B)$) coordinate (N)
			($(C)!1/3!(S)$) coordinate (P)
			($(D)!1/3!(S)$) coordinate (Q)
			;
			\foreach \i in{B,C,D}{\draw (S)--(\i);};
			\draw (B)--(C)--(D) (N)--(P)--(Q);
			\draw[dashed] (S)--(A)--(B) (A)--(D) (N)--(M)--(Q);
			\foreach \i/\g in {A/-90,B/-90,C/-90,D/-90,S/90,M/-60,N/-90,P/180,Q/60}
			\fill[black] (\i) circle(1pt)+(\g:3mm)node[scale=1]{$\i$};
		\end{tikzpicture}
	\end{center}
	\shortans{$6{,}2$}
	\loigiai{
		\begin{center}
			\begin{tikzpicture}[scale=0.5, font=\footnotesize,line join=round, line cap=round, >=stealth]
				\path 
				(0,0) coordinate (A)
				++(-140:3) coordinate (B)
				++(0:8) coordinate (C)
				($(A)+(C)-(B)$) coordinate (D)
				($(A)+(1,6)$) coordinate (S)
				($(D)!1/3!(A)$) coordinate (M)
				($(C)!1/3!(B)$) coordinate (N)
				($(C)!1/3!(S)$) coordinate (P)
				($(D)!1/3!(S)$) coordinate (Q)
				;
				\foreach \i in{B,C,D}{\draw (S)--(\i);};
				\draw (B)--(C)--(D) (N)--(P)--(Q);
				\draw[dashed] (S)--(A)--(B) (A)--(D) (N)--(M)--(Q);
				\foreach \i/\g in {A/-90,B/-90,C/-90,D/-90,S/90,M/-60,N/-90,P/180,Q/60}
				\fill[black] (\i) circle(1pt)+(\g:3mm)node[scale=1]{$\i$};
			\end{tikzpicture}
		\end{center}
		Ta có $MN$, $CD$, $PQ$ là ba giao tuyến của ba mặt phẳng $(ABCD)$, $(MNPQ)$, $(SCD)$.\\ Theo Định lý ba đường giao tuyến ta suy ra $MN \parallel CD \parallel PQ$.\\
		Theo tính chất song song, $SA \parallel (\alpha)$ nên $SA \parallel MQ$, với $MQ$ là giao tuyến của $(\alpha)$ với $(SAD)$.\\
		Theo Định lý Thales trong $\triangle SAD$ và $\triangle SCD$\\
		$\dfrac{PQ}{CD} = \dfrac{SQ}{SD} = \dfrac{AM}{AD} = \dfrac{2}{3} \Rightarrow PQ = \dfrac{2}{3} \cdot 9{,}3 = 6{,}2$.
	}
\end{ex}
\Closesolutionfile{ans}
\indapan{6}{ans/ans-CTST_1-OTGK1_De5-KQ}